% GENERATED, edit the French unit or the translation table
\mxunit{premiere-c10}{Trigonometry: circle, radians and related angles}\label{C10}
\begin{keypoint}[Chapter path]
\textbf{Objectives.} Convert from degrees to radians; place points; calculate standard sines and cosines; use the symmetries of the circle.

\textbf{Prerequisites.} Circle, Pythagoras, cosine and sine in a right triangle, symmetries, fractions.

\textbf{Yearly reference point.} 2--11 March: 8 hours\end{keypoint}
\section*{Diagnostic: \mxcount{diag} questions, 10 minutes}\label{C10-D}
Answer without consulting the lesson. Mark the questions for which help is needed. This diagnostic is not graded.
\begin{enumerate}
\item Circumference of a circle of radius 1?
\item How many degrees is a quarter turn?
\item Calculate $\sqrt {\dfrac{1}{2}}$.
\item What are the x- and y-coordinates of the reflection of $\pt{a;b}$ in the x-axis?
\item Simplify $\dfrac{5\pi}{3}-2\pi $.
\item In a right triangle, opposite side 3 and hypotenuse 5: sine?
\mxcounted{diag}
\end{enumerate}
\paragraph{After the diagnostic.} Revisit the flagged topics with the teacher; redo a similar question, then explain the rule used.
\paragraph{Speaking and logic.} Present a trigonometric identity involving related angles, naming the symmetry that justifies it; explain what a radian measures.
\mxnewpage
\section*{Investigate before formalising}\label{C10-RECH}
Construct, without a protractor, a point at angle $\dfrac{\pi}{3}$ then a point at angle $\dfrac{\pi}{4}$. Justify the exact values of their coordinates from the figures.
\paragraph{Organisation.} Individual investigation 5 min, comparison in pairs 5 min, class discussion 5 min. Keep a record of one attempt, even an unfinished one. A conjecture must then be justified.
\mxnewpage
\section{Measuring an angle in radians}\label{C10-U1}
\begin{keypoint}[Definition and results]
The unit circle is the circle with centre O and radius 1, oriented anticlockwise. One radian corresponds to an arc of length 1 on this circle. A full turn is $2\pi$ radians, a half turn $\pi$. On a circle of radius r, an arc length is $r\theta$ for a positive angle $\theta$ expressed in radians.
\end{keypoint}
\paragraph{Status and justification.} Definitions and proportionality between central angle and arc length. The conversion is $180^\circ=\pi$ radians; we multiply degrees by $\dfrac{\pi}{180}$.
\begin{example}$60^\circ =\dfrac{\pi}{3}$; $\dfrac{3\pi}{2}=270^\circ $. An angle of $\dfrac{\pi}{3}$ on a circle of radius 6 cm gives an arc of $2\pi$ cm.\end{example}\label{C10-EX1}
\begin{pitfall}The length of an arc is not that of the chord. In trigonometric calculations, check the degree or radian mode of the tool.\end{pitfall}
\begin{application}Convert $45^\circ$, $120^\circ$; calculate the arc length for radius 2 m and angle $\dfrac{3\pi}{4}$.\end{application}\label{C10-AP1}
\section{Wrapping and the image of a real number}\label{C10-U2}
\begin{keypoint}[Definition and results]
To each real number x corresponds a point M of the circle, obtained by an oriented path of length x starting from I(1;0). Two real numbers have the same image if their difference is an integer multiple of $2\pi$. A negative number indicates a path in the clockwise direction.
\end{keypoint}
\paragraph{Status and justification.} Adding or removing a whole number of turns keeps the end point. Conversely, returning to the same point requires a whole number of turns.
\begin{example}$\dfrac{13\pi}{3}=4\pi +\dfrac{\pi}{3}$ has the same image as $\dfrac{\pi}{3}$. $-\dfrac{\pi}{2}$ has the same image as $\dfrac{3\pi}{2}$: the point $\pt{0;-1}$.\end{example}\label{C10-EX2}
\begin{pitfall}Adding $\pi$ gives the opposite point, not the same point.\end{pitfall}
\begin{application}Place the images of $\dfrac{7\pi}{2}$ and $-\dfrac{5\pi}{3}$ by reducing them to within one turn.\end{application}\label{C10-AP2}
\section{Sine, cosine and standard values}\label{C10-U3}
\begin{keypoint}[Definition and results]
If M is the image of x, $M=\pt{\cos x;\sin x}$. So $-1\leq\cos x,\sin x\leq1$ and $\cos^2x+\sin^2x=1$. For $x=0,\dfrac{\pi}{6},\dfrac{\pi}{4},\dfrac{\pi}{3},\dfrac{\pi}{2}$, the cosines are $1,\dfrac{\sqrt 3}{2},\dfrac{\sqrt 2}{2},\dfrac{1}{2},0$ and the sines $0,\dfrac{1}{2},\dfrac{\sqrt 2}{2},\dfrac{\sqrt 3}{2},1$.
\end{keypoint}
\paragraph{Status and justification.} At $\dfrac{\pi}{4}$, the two coordinates are equal and positive; Pythagoras gives $2a^2=1$, so $a=\dfrac{\sqrt 2}{2}$. An equilateral triangle of side 1, cut by its altitude, has half-base $\dfrac{1}{2}$ and height $\dfrac{\sqrt 3}{2}$: $\cos (\dfrac{\pi}{3})=\dfrac{1}{2}$, $\sin (\dfrac{\pi}{3})=\dfrac{\sqrt 3}{2}$. The complementary angle gives $\dfrac{\pi}{6}$.
\begin{example}The point at angle $\dfrac{\pi}{3}$ is $\pt{\dfrac{1}{2};\dfrac{\sqrt 3}{2}}$. The scalar product of two unit vectors forming this angle equals its cosine, so $\dfrac{1}{2}$. The geometry chapters thus connect to the circle.\end{example}\label{C10-EX3}
\mxinclude{figures/C10-U3}
\begin{pitfall}The cosine is the x-coordinate, the sine the y-coordinate. The identity gives an absolute value: the quadrant determines the sign.\end{pitfall}
\begin{application}If $\cos x=\dfrac{3}{5}$ and the image of x is in the fourth quadrant, give $\sin x$.\end{application}\label{C10-AP3}
\section{Related angles and symmetries}\label{C10-U4}
\begin{keypoint}[Definition and results]
The symmetries of the circle give $\cos(-x)=\cos x$, $\sin(-x)=-\sin x$; $\cos(\pi-x)=-\cos x$, $\sin(\pi-x)=\sin x$; $\cos(\pi+x)=-\cos x$, $\sin(\pi+x)=-\sin x$. Adding $2k\pi$, k an integer, changes neither sine nor cosine.
\end{keypoint}
\paragraph{Status and justification.} Each equality can be read from the coordinates of the symmetric points: horizontal axis for $-x$, vertical axis for $\pi-x$, centre O for $\pi+x$. This is not a catalogue without justification.
\begin{example}$\cos (\dfrac{5\pi}{6})=-\dfrac{\sqrt 3}{2}$, $\sin (\dfrac{5\pi}{6})=\dfrac{1}{2}$ by symmetry from $\dfrac{\pi}{6}$. $\cos (-\dfrac{\pi}{4})=\dfrac{\sqrt 2}{2}$, $\sin (-\dfrac{\pi}{4})=-\dfrac{\sqrt 2}{2}$.\end{example}\label{C10-EX4}
\begin{pitfall}Two angles with the same cosine do not necessarily have the same sine. Reduce the angle, then read the signs.\end{pitfall}
\begin{application}Calculate the two coordinates of the images of $\dfrac{4\pi}{3}$ and $\dfrac{7\pi}{4}$.\end{application}\label{C10-AP4}
\section{Approximating pi and putting it to use}\label{C10-U5}
\begin{keypoint}[Definition and results]
A regular inscribed polygon with n sides in the unit circle has half-perimeter $n\sin (\dfrac{\pi}{n})$; a circumscribed polygon has half-perimeter $n\tan (\dfrac{\pi}{n})$, with $\tan \theta =\dfrac{\sin \theta}{\cos}\theta $ here. Archimedes' method doubles the number of sides and calculates the new lengths by Pythagoras, without knowing pi in advance.
\end{keypoint}
\paragraph{Status and justification.} The inscribed perimeter is less than that of the circle and the circumscribed one greater: these geometric bounds are accepted. If c is the chord of one side, let $h=\sqrt {1-\dfrac{c^2}{4}}$; the chord of the half-angle is $\dfrac{c}{\sqrt {2(1+h)}}$. This formula follows from Pythagoras and avoids subtracting nearly equal numbers.
\begin{example}For the inscribed hexagon, c=1 and half-perimeter 3. The circumscribed one gives $2\sqrt3\approx\num{3,4641}$. By doubling the sides, we tighten the bounds. The supplied program starts from n=6 and c=1.\end{example}\label{C10-EX5}
\begin{pitfall}Using a sine function directly with pi already known would check a calculation, but would not construct an approximation independent of pi.\end{pitfall}
\begin{application}For n=6 and c=1, calculate the new chord after doubling.\end{application}\label{C10-AP5}
\mxnewpage
\section*{Exercise path}\label{C10-PARC}
Start with exercises 1 and 2. Then choose one of exercises 3 or 4 and do it without help. Exercises 5 to 7 are a reserve for transfer, investigation or programming: the teacher designates one. Exercise 8 is revisited after a delay.
\begingroup
\mxexercisehead{C10-E1}{Converting and measuring}
We use the unit circle to measure angles in radians.
\begin{enumerate}[label=\arabic*.]
\item Convert $150^\circ$ and $-90^\circ$ to radians, showing the conversion factor.
\item A circular sector has radius 4 cm and an angle of $\dfrac\pi3$ radians. Calculate the exact length of its arc.
\item Explain why an angle measure can be negative while a geometric arc length is positive.
\end{enumerate}
\mxexercisehead{C10-E2}{Reducing angles}
Consider the real numbers $\dfrac{17\pi}{6}$, $-\dfrac{7\pi}{4}$ and $5\pi$.
\begin{enumerate}[label=\arabic*.]
\item For each, determine a representative in $\intfo{0}{2\pi}$ having the same image on the unit circle. Show the number of turns added or removed.
\item Place the three images on a circle and identify them by the starting real numbers.
\end{enumerate}
\mxexercisehead{C10-E3}{Exact values}
We look for the exact values of the cosine and sine of $\dfrac{2\pi}{3}$, $-\dfrac\pi6$ and $\dfrac{5\pi}{4}$.
\begin{enumerate}[label=\arabic*.]
\item Place the three images on the unit circle. For each, state the expected sign of its two coordinates.
\item Calculate the exact values using a standard angle and a symmetry. Present the answer in a table: angle, cosine, sine.
\end{enumerate}
\mxexercisehead{C10-E4}{Recovering the sign}
A real number $x$ satisfies $\sin x=\dfrac5{13}$. Its image on the unit circle lies in the second quadrant, that is to the left of the y-axis and above the x-axis.
\begin{enumerate}[label=\arabic*.]
\item Determine the exact value of $\cos x$. Justify both its absolute value and its sign.
\item Explain why the value of $\sin x$ alone would not have been enough to choose the sign of $\cos x$.
\end{enumerate}
\mxexercisehead{C10-E5}{A wheel}
A wheel of radius $\num{0,30}$ m rolls on horizontal ground without slipping. It turns through an angle of $\dfrac{5\pi}{3}$ radians. We look for the distance travelled by its axle in the direction of motion.
\begin{enumerate}[label=\arabic*.]
\item Calculate this distance in exact form, explaining the link with the arc length unrolled.
\item Explain why the length of a chord of the circle must not be used. Is the path of a point fixed on the rim of the wheel, seen from the ground, the same as that of the axle?
\end{enumerate}
\mxexercisehead{C10-E6}{Chord and scalar product}
On the unit circle with centre $O$, the point $A$ is the image of 0 and the point $B$ that of $\dfrac{2\pi}{3}$.
\begin{enumerate}[label=\arabic*.]
\item Determine the exact coordinates of $A$ and $B$, then calculate the length $AB$.
\item Recover this length by applying the Al-Kashi relation to the triangle $OAB$. State the lengths and the angle used.
\end{enumerate}
\mxexercisehead{C10-E7}{Reading an equation on the circle}
We look for the real numbers $x$ in the interval $\intfo{0}{2\pi}$ whose image on the unit circle has x-coordinate $\dfrac12$.
\begin{enumerate}[label=\arabic*.]
\item Translate the condition on the x-coordinate into an equation involving $\cos x$.
\item Draw the corresponding vertical line on a circle, mark its intersections and give the two solutions. Check that they belong to the interval required.
\end{enumerate}
\mxexercisehead{C10-E8}{Delayed review}
Consider the angle of measure $\dfrac{25\pi}{4}$ radians. The answers must be exact, with no decimal approximation.
\begin{enumerate}[label=\arabic*.]
\item Reduce this angle, showing a whole number of complete turns.
\item Deduce its cosine and sine. Justify the values using the unit circle.
\end{enumerate}
\paragraph{Homework.} Do only the task designated at the end of the session, in the dated progression: 10--20 min at most. The exercises depend on the topics already studied, so some questions are deferred. Revisit exercise 8 one week after the chapter (10 min). If the allotted time is exceeded, note the difficulty and bring your attempts; do not pile up the whole reserve.
\endgroup
\mxneed{12}\section*{To do in the CQFD textbook}
These exercises are done in your textbook. The progression indicates the session; they replace the homework planned for that day (10 to 20 minutes in total).
\paragraph{CQFD, exercise 32 p.~220, questions a and b.}\label{C10-M1} Use the exact value supplied and the related angles, without a calculator.
\paragraph{CQFD, exercise 39 p.~220, questions 1, 2, 3 and 5.}\label{C10-M2} Do the four calculation and circle-reading questions; justify the answers on your paper.
\begingroup\small
\mxnewpage\section*{MCQ 1: understand and choose}\label{C10-Q1}
\mxCount{qcm1} single-answer questions for practice. Choose one answer per question; then explain one distractor. This chapter format does not reproduce the official EAM format.
\begin{enumerate}
\item $180^\circ$ equals
\begin{enumerate}[label=\Alph*.,itemsep=2pt]
\item $2\pi$
\item $\pi$
\item $\dfrac{\pi}{2}$
\item $1$
\end{enumerate}
\item $\cos (\dfrac{\pi}{3})=$
\begin{enumerate}[label=\Alph*.,itemsep=2pt]
\item $\dfrac{1}{2}$
\item $\dfrac{\sqrt 3}{2}$
\item $0$
\item $1$
\end{enumerate}
\item $\sin (-\dfrac{\pi}{2})=$
\begin{enumerate}[label=\Alph*.,itemsep=2pt]
\item $1$
\item $0$
\item $\dfrac{1}{2}$
\item $-1$
\end{enumerate}
\item $\cos^2x+\sin^2x=$
\begin{enumerate}[label=\Alph*.,itemsep=2pt]
\item $0$
\item $2$
\item $1$
\item $x$
\end{enumerate}
\item The same point corresponds to angles differing by
\begin{enumerate}[label=\Alph*.,itemsep=2pt]
\item $\pi$
\item $2k\pi$, k an integer
\item $\dfrac{\pi}{2}$
\item $1$
\end{enumerate}
\mxcounted{qcm1}
\end{enumerate}
\endgroup
\begingroup\small
\mxneed{18}\section*{MCQ 2: understand and choose}\label{C10-Q2}
\mxCount{qcm2} single-answer questions for practice. Choose one answer per question; then explain one distractor. This chapter format does not reproduce the official EAM format.
\begin{enumerate}
\item $\cos(\pi-x)=$
\begin{enumerate}[label=\Alph*.,itemsep=2pt]
\item $\cos x$
\item $\sin x$
\item $-\cos x$
\item $-\sin x$
\end{enumerate}
\item $\sin (\dfrac{\pi}{4})=$
\begin{enumerate}[label=\Alph*.,itemsep=2pt]
\item $\dfrac{\sqrt 2}{2}$
\item $\dfrac{1}{2}$
\item $\dfrac{\sqrt 3}{2}$
\item $\sqrt2$
\end{enumerate}
\item Image of $\dfrac{3\pi}{2}$:
\begin{enumerate}[label=\Alph*.,itemsep=2pt]
\item $\pt{1;0}$
\item $\pt{-1;0}$
\item $\pt{0;1}$
\item $\pt{0;-1}$
\end{enumerate}
\item Arc of radius 3 and angle 2 radians: length
\begin{enumerate}[label=\Alph*.,itemsep=2pt]
\item $\dfrac{2}{3}$
\item $6$
\item $\dfrac{3}{2}$
\item $6\pi$
\end{enumerate}
\item Interior of the second quadrant: sign of the cosine
\begin{enumerate}[label=\Alph*.,itemsep=2pt]
\item positive
\item zero everywhere
\item negative
\item cannot be determined
\end{enumerate}
\mxcounted{qcm2}
\end{enumerate}
\subsection*{True or false? Justify.}\label{C10-VF}
\begin{enumerate}
\item $\cos(-x)=-\cos x$ for all x.
\item $\sin(x+2\pi)=\sin x$.
\item The value of $\sin x$ always determines $\cos x$.
\end{enumerate}
\begin{keypoint}[Self-assessment]I reduce angles by complete turns; I read cosine and sine in the right order; I justify the signs; I keep the exact values.

For each objective, indicate: achieved without help; achieved with help; to be reworked. Write one example and one mistake to avoid.\end{keypoint}
\endgroup
\ifmxteacher\mxinclude{corriges/C10}\fi
